\section*{Chapitre \ref{ch:b3:planetary-atmospheres} --- Atmosphères planétaires et habitabilité}

\begin{solution}{exo:b3:planetary-atmospheres:1}
$H = RT/Mg$~: Vénus $8.314\times737/(0.04401\times8.87) = \qty{15.7}{km}$~;
la Terre $8.314\times288/(0.02897\times9.80) = \qty{8.4}{km}$~; Mars
$8.314\times214/(0.04334\times3.71) = \qty{11.1}{km}$; Titan
$8.314\times94/(0.0280\times1.35) = \qty{20.7}{km}$. L'atmosphère de Titan,
malgré sa froideur, est la plus étendue~: sa gravité est sept fois plus faible
que celle de la Terre.
\end{solution}

\begin{solution}{exo:b3:planetary-atmospheres:2}
$\lambda = GMm/kTr = \num{6.674e-11}\times\num{5.97e24}\times\num{1.674e-27}
/(\num{1.381e-23}\,T\times\num{6.878e6})$~: 8.78 à \qty{800}{K}, 4.39 à
\qty{1600}{K}. Les fractions d'hydrogène qui s'échappent sont $9.78\,e^{-8.78} =
\num{1.5e-3}$ et $5.39\,e^{-4.39} = \num{6.7e-2}$~: le cycle solaire
modifie la fraction d'hydrogène qui s'échappe d'un facteur d'environ 44.
\end{solution}

\begin{solution}{exo:b3:planetary-atmospheres:3}
(a)~$d = S_\mathrm{eff}^{-1/2}$~: \omterm{def:b3:planetary-atmospheres:runaway}{effet de serre humide} \qty{0.993}{au}, \omterm{def:b3:planetary-atmospheres:runaway}{effet de serre galopant}
\qty{0.975}{au}, effet de serre maximal \qty{1.705}{au}. (b)~Vénus, à 1.91,
reçoit davantage que le flux \omterm{def:b3:planetary-atmospheres:runaway}{limite de l'effet de serre galopant} et même davantage que la limite empirique
de Vénus récente (1.78)~: elle se trouve à l'intérieur du bord interne. Mars, à 0.43,
se trouve entre 1.014 et 0.3438, à l'intérieur de la zone habitable conservative~:
sa sécheresse vient de sa petite taille et de la perte de son atmosphère, et non
de sa distance.
\end{solution}

\begin{solution}{exo:b3:planetary-atmospheres:4}
(a)~$\log_{10}p = 3.9895 - 443.028/93.66 = -0.741$, $p =
\qty{0.18}{bar}$. (b)~$0.05\times1.6 = \qty{0.080}{bar}$, une humidité relative
de $0.080/0.18 = 44\,\%$. (c)~L'air au-dessus des mers est loin
d'être saturé~: le méthane s'évapore de celles-ci, comme l'eau le fait des
océans terrestres, et l'air humide doit s'élever et se refroidir avant qu'il ne pleuve.
\end{solution}

\begin{solution}{exo:b3:planetary-atmospheres:5}
$\lambda_\mathrm D = 7.02\times2.014/1.008 = 14.0$. $u_\mathrm D/u_\mathrm H
= \sqrt{1.008/2.014} = 0.71$. Rapport par atome $= 0.71\times(15.0\,
e^{-14.0})/(8.02\,e^{-7.02}) = 0.71\times\num{1.69e-3} = \num{1.2e-3}$.
L'\omterm{def:b3:planetary-atmospheres:jeans-escape}{échappement de Jeans} laisse presque aucun deutérium s'échapper~: c'est un filtre puissant, avec
$\alpha \approx 0.001$.
\end{solution}

\begin{solution}{exo:b3:planetary-atmospheres:6}
(a)~$0.016/\num{1.5575e-4} = 103$. (b)~Avec $\alpha = 0$, l'enrichissement est
$M_0/M$~: $M_0 = 103\times\num{6e15} = \qty{6.2e17}{kg}$, soit
$\num{6.2e17}/(1000\times\num{4.60e14}) = \qty{1.3}{m}$ d'eau sur
Vénus. (c)~$M_0/M = 103^{1/0.2} = \num{1.2e10}$, un inventaire passé de
\qty{7e25}{kg}, soit douze fois la masse de la Terre~: impossible. Le
enrichissement exige un échappement fortement fractionnant (de Jeans ou non thermique),
actif au moins pendant la perte des dernières traces d'eau, ou un
état stationnaire avec un apport d'eau.
\end{solution}

\begin{solution}{exo:b3:planetary-atmospheres:7}
(a)~En quatre milliards d'années ($\num{1.26e17}$~s), Mars perd
$\num{3.8e43}$ atomes, $\num{1.9e43}$ molécules d'eau, $\qty{5.7e17}{kg}$~;
sur $4\pi(\num{3.3895e6})^2 = \qty{1.44e14}{m^2}$, soit une couche
\qty{3.9}{m} d'épaisseur. (b)~Vingt fois plus, \qty{78}{m}. (c)~Même le
taux de forte obliquité n'atteint que le bas de la fourchette de 100 à
\qty{1500}{m}~; le reste s'est enfoui dans la croûte et la glace, ou a été perdu
plus vite lorsque le Soleil était jeune.
\end{solution}

\begin{solution}{exo:b3:planetary-atmospheres:8}
(a)~$m = pA/g = 636\times\num{1.44e14}/3.71 = \qty{2.5e16}{kg}$.
(b)~$0.1\times\num{1.26e17} = \qty{1.3e16}{kg}$, soit la moitié
de l'atmosphère actuelle~; la totalité en $\num{2.5e16}/0.1 = \qty{2.5e17}{s}$,
\qty{7.8}{Ga}. (c)~Le jeune Soleil était plus actif et ses
tempêtes plus fréquentes, et la perte croît pendant les tempêtes~; d'autres voies
(hydrogène, pulvérisation, échappement photochimique) ne sont pas pris en compte.
\end{solution}

\begin{solution}{exo:b3:planetary-atmospheres:9}
(a)~$S(1 - 0.193)/4 = 291$ donne $S = \qty{1442}{W.m^{-2}}$, soit 1.06 fois
celle de la Terre. (b)~Aujourd'hui, Vénus absorberait $0.807\times2601/4 =
\qty{525}{W.m^{-2}}$~; une Vénus jeune, $0.7\times2601 = \qty{1821}{W.m^{-2}}$
reçu, \qty{367}{W.m^{-2}} absorbé~: les deux dépassent la limite, donc une
Vénus jeune sans nuages aurait basculé dès le départ dans l'\omterm{def:b3:planetary-atmospheres:runaway}{effet de serre galopant}. (c)~Les nuages
augmentent l'albédo~; des modèles tridimensionnels incluant les nuages ont trouvé une voie
vers une Vénus jeune habitable.
\end{solution}

\begin{solution}{exo:b3:planetary-atmospheres:10}
(a)~$H = kT/mg = \num{1.381e-23}\times1000/(16\times\num{1.661e-27}\times9.21)
= \qty{56}{km}$. (b)~$n_\mathrm{exo} = 1/(\sigma_c H) = 1/(\num{3e-19}\times
\num{5.64e4}) = \qty{5.9e13}{m^{-3}}$; $z = 200 + 56.4\ln(\num{1e16}/
\num{5.9e13}) = 200 + 56.4\times5.13 = \qty{490}{km}$. (c)~Une gravité plus faible
en altitude fait croître la hauteur d'échelle avec l'altitude, de sorte que la densité
décroît plus lentement que prévu~: l'\omterm{def:b3:planetary-atmospheres:exobase}{exobase} s'élève.
\end{solution}

\begin{solution}{exo:b3:planetary-atmospheres:11}
(a)~$\Phi = \num{3.0e13}\times\num{3e-3} = \qty{9e10}{cm^{-2}.s^{-1}} =
\qty{9e14}{m^{-2}.s^{-1}}$; sur $4\pi(\num{6.371e6})^2 =
\qty{5.10e14}{m^2}$, $\num{4.6e29}$ molécules de $\ce{H2}$ par seconde.
(b)~Un océan contient $\num{1.4e24}/18.015\times\num{6.022e23} =
\num{4.7e46}$ molécules d'eau, à raison d'une molécule de $\ce{H2}$ chacune~: cela dure
$\qty{1.0e17}{s} = \qty{3.2}{Ga}$, comparable à l'âge de
la Terre, ce qui explique pourquoi un rapport de mélange de \num{3e-3} marque la
limite de l'\omterm{def:b3:planetary-atmospheres:runaway}{effet de serre humide}.
\end{solution}

\begin{solution}{exo:b3:planetary-atmospheres:12}
(a)~$p_1 = g/D\kappa = 9.81/(1.66\times0.005) = \qty{1180}{Pa}$.
(b)~$1/T = 1/273.16 - \ln(1182/611.657)\times461.5/\num{2.5e6}$, $T =
\qty{282.5}{K}$, et $\sigma T^4 = \qty{361}{W.m^{-2}}$. (c)~La limite réelle
est plus basse, donc le niveau de rayonnement est plus froid et à plus basse pression~: la
vapeur d'eau réelle est plus opaque que ce gaz gris, et son $\kappa$ effectif
plus grand.
\end{solution}

\begin{solution}{pb:b3:planetary-atmospheres:1}
\textbf{1.} $g = GM/R^2 = \num{6.674e-11}\times\num{5.2e24}/(\num{6.2e6})^2 =
\qty{9.03}{m.s^{-2}}$; $\bar\rho = M/\tfrac43\pi R^3 = \qty{5210}{kg.m^{-3}}$.

\textbf{2.} $v_\mathrm{esc} = \sqrt{2GM/r} = \sqrt{2\times\num{3.47e14}/\num{6.7e6}}
= \qty{10.18}{km.s^{-1}}$.

\textbf{3.} $u = \sqrt{2kT/m}$~: hydrogène \qty{4.81}{km.s^{-1}}, hélium
\qty{2.41}{km.s^{-1}}, oxygène \qty{1.21}{km.s^{-1}}.

\textbf{4.} $\lambda = (v_\mathrm{esc}/u)^2$: 4.49, 17.8, 71.2.

\textbf{5.} Seul l'hydrogène~: sa fraction échappée vaut 0.06, celle de l'hélium
\num{3.5e-7}, celle de l'oxygène \num{9e-30}. Aucun n'est proche du régime hydrodynamique,
$\lambda \approx 1.5$ à 3.

\textbf{6.} $\lambda = 4.49\times1400/3000 = 2.09$~: dans la zone de transition~;
l'hydrogène se serait échappé en un flux hydrodynamique, entraînant peut-être
des gaz plus lourds.

\textbf{7.} $(1 + 4.49)\,e^{-4.49} = 0.0618$.

\textbf{8.} $n u/2\sqrt\pi = \num{2e11}\times4806/3.545 = \qty{2.71e14}{m^{-2}.s^{-1}}$;
$\Phi_\mathrm J = \num{2.71e14}\times0.0618 = \qty{1.68e13}{m^{-2}.s^{-1}}$.

\textbf{9.} $4\pi(\num{6.7e6})^2 = \qty{5.64e14}{m^2}$~; $\num{9.5e27}$ atomes par
seconde.

\textbf{10.} $\num{1.4e24}/18.015 = \qty{7.8e22}{mol}$, $\num{4.68e46}$ molécules,
$\num{9.36e46}$ atomes d'hydrogène.

\textbf{11.} $\num{9.36e46}/\num{9.5e27} = \qty{9.9e18}{s}$, soit environ 310 milliards
d'années.

\textbf{12.} $\lambda_\mathrm D = 4.49\times2.014/1.008 = 8.96$~; fraction échappée
$9.96\,e^{-8.96} = \num{1.28e-3}$~; efficacité relative $0.707\times
\num{1.28e-3}/0.0618 = 0.015$.

\textbf{13.} $H_a = kT/m_a g = \num{1.381e-23}\times200/(44.01\times\num{1.661e-27}
\times9.03) = \qty{4.18}{km}$.

\textbf{14.} $b = \qty{1.2e21}{m^{-1}.s^{-1}}$, donc
\[
\Phi_\mathrm{DL} = \frac{b f(m_a - m_i)}{m_a H_a}
= \frac{\num{1.2e21}\times\num{1e-3}\times(42.0/44.0)}{4185}
= \qty{2.74e14}{m^{-2}.s^{-1}}.
\]

\textbf{15.} $\num{3.0e13}\times\num{1e-3}\times9.03/9.81 =
\qty{2.76e10}{cm^{-2}.s^{-1}} = \qty{2.76e14}{m^{-2}.s^{-1}}$~: le même
à 1\,\% près.

\textbf{16.} $4\pi(\num{6.2e6})^2 = \qty{4.83e14}{m^2}$~; $\num{1.32e29}$ molécules
de $\ce{H2}$, $\num{2.64e29}$ atomes par seconde.

\textbf{17.} La limite de diffusion fournit 28 fois ce que l'\omterm{def:b3:planetary-atmospheres:exobase}{exobase} laisse échapper à
sa densité actuelle. L'hydrogène s'accumule dans la thermosphère jusqu'à ce que le
flux de Jeans égale l'apport~: à long terme, l'échappement est
limité par la diffusion, comme sur la Terre.

\textbf{18.} $\num{2e11}\times28 = \qty{5.6e12}{m^{-3}}$.

\textbf{19.} $\num{4.68e46}/\num{1.32e29} = \qty{3.54e17}{s}$, 11.2 milliards d'années.

\textbf{20.} $S = 0.30/0.60^2 = 0.833$.

\textbf{21.} $d = \sqrt{L/S_\mathrm{eff}}$~: bord intérieur $\sqrt{0.30/0.939} =
\qty{0.565}{au}$, bord extérieur $\sqrt{0.30/0.290} = \qty{1.02}{au}$. Kepos,
à \qty{0.60}{au}, se trouve à l'intérieur, près du bord intérieur.

\textbf{22.} $0.939/0.833 = 1.13$~: une hausse de 13\,\% du flux.

\textbf{23.} $0.833\times1361\times0.7/4 = \qty{198}{W.m^{-2}}$, $T_\mathrm{eq} =
(198/\sigma)^{1/4} = \qty{243}{K}$~; bien en dessous de la
\qty{291}{W.m^{-2}} limite (calculée pour une étoile de type solaire, mais la marge
est large).

\textbf{24.} Trois fois plus rapide~: 3.7 milliards d'années.

\textbf{25.} \textbf{11~milliards~d'années}.
\end{solution}
